\section{Tail counts and composition at a common maximum}\label{sec:cumulativecores}

We distinguish log-concavity of a coefficient sequence from log-concavity
of its tail sums. For a nonnegative polynomial $P$ of degree $d$, write
\[
 \mathcal T_P(k)=\sum_{j=k}^{d}P_j\qquad(k\in\mathbb Z),
\]
where coefficients outside $[0,d]$ vanish. Thus $\mathcal T_P$ is constant
for $k\le0$ and zero for $k>d$.

\begin{lemma}[Cumulative convolution closure]
Let $p,q$ be nonzero nonnegative finitely supported sequences. If their
cumulative sums are log-concave on all integer indices, the cumulative
sums of $p*q$ are log-concave as well.
\end{lemma}
\begin{proof}
Normalize $p$ and $q$ by their positive total masses and translate their
supports into the nonnegative integers. The hypothesis says precisely
that both distributions are discrete decreasing reversed hazard
distributions. Apply Nanda and Sengupta~\cite{NandaSengupta},
Theorem 4.2(d). Undo the normalization and the translations.
\end{proof}

\begin{lemma}[Closure at a common maximum]\label{lem:cumulativegraft}
Fix $M\ge2$ and arbitrary hosts $X,Y$ with parts at most $M$.
Suppose the cumulative sums of each of $q_X,q_Y,\delta_X,\delta_Y$
are log-concave, where the two kernels use peak $M$.
Then both the reversed rank polynomial and the reversed peak-$M$
kernel of $XMY$ have log-concave cumulative sums.
\end{lemma}
\begin{proof}
The maximum identity gives
\[
 D_{XMY}=D_XD_Y,\qquad
 q_{XMY}=q_X\delta_Y+z^{a+M}q_Y,
 \qquad a=\wt X,\quad b=\wt Y.
\]
The first identity follows by expanding
$1+(S_M-1)(F_X+F_Y+(S_M-1)F_XF_Y)$.
After reversing, it gives $\delta_{XMY}=\delta_X\delta_Y$;
cumulative convolution closure proves the kernel assertion.

For the rank polynomial, interchange the arms if necessary so that
$a\ge b$. This leaves the polynomial unchanged. The cumulative sums of
$q_X\delta_Y$ are log-concave. At centers $r\le a+M-2$, the
added term has not entered any of the three cumulative values.
At every remaining center, the consecutive increments of the full
cumulative sequence satisfy
\[
 (q_{XMY})_r\ge(q_{XMY})_{r+1}.
\]
Indeed, their forward ranks are $a+b+M-r\le b+1$ and the preceding
rank, while $b+1\le\lceil(a+b+M)/2\rceil$. Apply
Lemma~\ref{lem:half}; ranks outside the support cause no difficulty.
Nonnegative decreasing increments imply concavity and hence
log-concavity of the cumulative sequence. This includes the joining
center $a+M-1$.
\end{proof}

\begin{theorem}[Valleys separated by a common maximum]
\label{thm:separated}
Let
\[
 w=Z_0\,M\,Z_1\,M\cdots M\,Z_s,
\]
where each $Z_i$ weakly decreases and then weakly increases, possibly
with one or both segments empty, and every part of $Z_i$ is strictly
below $M$.
There is no restriction on $M$, $s$, or the block lengths.
Both $F_w$ and $1+(S_M-1)F_w$ have log-concave tail counts.
In particular, the number of descents need not be bounded.
\end{theorem}
\begin{proof}
For $M=1$, all blocks are empty, and the word consists of ones;
both conclusions follow directly from the corresponding intervals.
Assume $M\ge2$. Every nonempty valley block has a global maximum at
an endpoint; removing that maximum leaves a valley block. Word reversal
preserves the rank polynomial, so removing either endpoint maximum $t$
uses the identity $F=1+S_tF_{\mathrm{rest}}$. Induction shows that a
valley block has the same rank polynomial as its increasing rearrangement.
Thus each block has the two cumulative properties proved for monotone
arms in the preceding section. Repeated application of
Lemma~\ref{lem:cumulativegraft} proves both conclusions. Equivalently,
the kernel factorization is
\[
 D_w=\prod_{i=0}^s D_{Z_i}.
\]
This factorization concerns the kernels at the same peak $M$.
\end{proof}

For example, $(1,7,10,8,3,10,2,9,10,7,4,1)$ has five descents and
satisfies the theorem with $M=10$. The theorem remains a statement about
cumulative counts; it makes no rank-unimodality assertion for this
larger class. Appending or prepending successive global maxima also
preserves the rank-tail conclusion, by the reverse append argument.

\begin{theorem}[Log-concavity of one-descent tail counts]
Let $w$ be any composition with at most one adjacent descent, with no
bound on its length or part sizes. Then
\[
 \mathcal T_{F_w}(k)^2\ge
 \mathcal T_{F_w}(k-1)\mathcal T_{F_w}(k+1)
 \qquad(k\in\mathbb Z).
\]
The same conclusion holds for every bitonic composition. For increasing
strict arms $X,Y$, it also holds for the relative polynomial $F_XD_Y$.
\end{theorem}

\begin{proof}
Consider first a strict core $(X,M,Y)$. The proved cumulative reverse-arm and reverse-kernel
theorems, together with cumulative convolution closure, show that the
cumulative sums of $q_X\delta_Y$ are log-concave. This proves the
relative-polynomial assertion, by reversing the index.
The core itself is covered by Theorem~\ref{thm:separated} with two
blocks.

Monotone hosts are already covered by the reverse-arm theorem.
Appending a global maximum changes $F$ to $1+S_tF$ and preserves
log-concavity of the cumulative reverse profile by the same window and
terminal-step argument used for monotone arms. That argument requires
only a nonnegative polynomial with constant coefficient one, a
log-concave cumulative reverse profile, and the maximum append identity.
Strip successive global maxima from the terminal increasing segment of
a one-descent word until the strict core is reached, and then restore
them. Inserting another copy next to the central peak also changes
$F$ to $1+S_MF$: its peak-$M$ kernel is multiplied by $S_M$, by the
maximum identity. Thus repeated copies of the peak are restored by
the same cumulative argument.
Finally, reversing the decreasing arm of a bitonic word does not change
its arm polynomial; the maximum decomposition gives the identical
core polynomial. This proves all assertions.
\end{proof}

The theorem asserts log-concavity of tail \emph{counts}, not tail
log-concavity of the individual rank coefficients. The latter is a
different, stronger shape condition. Log-concavity of tail sums alone
does not imply rank-unimodality.

\paragraph{External input and verification boundary.}
The cumulative convolution lemma is an external theorem. Ordinary
convolution closure for log-concave coefficient sequences does not
justify this step. The supplied Lean files verify the reverse append
and escape counting models; they do not formalize this external
probability theorem or the complete core theorem.
